\section{Périmètre et organisation du travail}
UrbanFlow analyse les enregistrements de trajets des taxis jaunes de New York.
Le jeu initial couvre janvier, février et mars 2024, dont le changement d'heure
du 10 mars. Avril est réservé à la mise à jour F6 ; mai et juin complètent les
six mois prévus pour les expériences. Brooklyn sert de point d'appui aux exemples,
mais les données de toute la ville sont conservées. Les trajets enregistrés ne
permettent pas de mesurer une demande non satisfaite ni l'ensemble des déplacements
new-yorkais : les résultats portent sur cette archive de taxis \cite{sujet}.

Nous nous sommes réparti le travail : Théo s'est occupé de la récupération des données, de leur qualité, du modèle, de l'API locale et des services F1 à F5, et Ion de F6, des déploiements répartis, des expériences E1-E3 et de leur analyse. Pour la partie données et F1-F5, 
Codex a été utilisé pour développer et expliquer du code, préparer des vérifications et
rédiger. 



\section{Données et règles de qualité}\label{sec:donnees}
Les sources sont les Parquet mensuels \emph{Yellow Taxi Trip Records} et le
référentiel officiel des zones TLC. Le manifeste donne pour chaque fichier son
URL, sa date de récupération, sa taille et son SHA-256. Nous avons récupéré les fichiers de janvier à avril le 2 novembre 2026 et ceux de mai et juin le 3 novembre ;
les empreintes complètes sont dans \chemin{data/manifest.json} \cite{tlc,preuves}.

\begin{center}\small
\begin{tabular}{lrrrr}
\toprule
Fichier 2024 & Lignes sources & Profil activity & Profil flows & Profil journey\\
\midrule
Janvier & 2\,964\,624 & 2\,929\,120 & 2\,898\,055 & 2\,841\,251\\
Février & 3\,007\,526 & 2\,971\,445 & 2\,942\,952 & 2\,878\,111\\
Mars & 3\,582\,628 & 3\,538\,317 & 3\,503\,542 & 3\,420\,628\\
Avril & 3\,514\,289 & 3\,470\,540 & 3\,435\,410 & 3\,394\,040\\
Mai & 3\,723\,833 & 3\,674\,841 & 3\,636\,771 & 3\,591\,918\\
Juin & 3\,539\,193 & 3\,490\,398 & 3\,452\,592 & 3\,404\,607\\
\midrule
Total & 20\,332\,093 & 20\,074\,661 & 19\,869\,322 & 19\,530\,555\\
\bottomrule
\end{tabular}
\end{center}
Ces effectifs sont ceux des bilans par fichier : 9\,554\,778 lignes sources pour les
trois mois initiaux, 20\,332\,093 pour les six mois, pas un nombre garanti de
déplacements physiques distincts. À l'import, aucune ligne d'un fichier n'a été
retrouvée dans un autre : la base contient autant de documents que de lignes
sources. Le référentiel contient 265 zones. Dans le fichier de janvier,
25\,258 lignes ont une zone de départ rattachée à Brooklyn ; Manhattan en compte
2\,646\,948 \cite{preuves}.

Une ligne appartient à une période selon sa date réelle de départ, et non selon
le mois du fichier. F1 et F3 en mode arrivées utilisent toutefois la zone \emph{et}
la date d'arrivée. Les intervalles sont semi-ouverts $[d_0,d_1[$ en heure locale
\code{America/New_York}. Cette convention évite de compter deux fois une borne.

Les anomalies de qualité ne provoquent pas de suppression à l'import : elles sont conservées dans
\code{quality_flags}, puis chaque analyse applique son profil.
\code{activity} exclut les montants totaux négatifs et les doublons exacts ;
\code{flows} ajoute les zones 264/265, sans localisation précise ;
\code{journey} ajoute les durées nulles ou négatives, celles au-delà de six heures,
les distances nulles ou supérieures à 100 miles, les vitesses supérieures à
80 miles/h et les heures DST non résolues. Le profil activity sert à F1 et F3, flows à F2 et F5, journey à F4. Ces seuils larges écartent des anomalies
manifestes, sans définir un trajet « normal ». En janvier, le profil journey
écarte 123\,373 lignes, soit 4,16\,\%, contre 1,20\,\% pour activity.
Les motifs peuvent se cumuler ; leur somme n'est pas le total des exclusions.

Un doublon exact est une copie de toutes les colonnes dans un même fichier ;
les deux lignes restent stockées et seule la première entre dans les analyses. Une annulation est repérée par
le fournisseur, les deux horodatages, les zones et un montant opposé ; l'original
positif est conservé. La détection est faite sur le fichier entier avant l'écriture
par lots de 5\,000, ce qui évite de dépendre des frontières de lots, mais impose
un coût mémoire que nous n'avons pas mesuré. Les anomalies ne sont pas toutes
couvertes : le contrôle des zones vise notamment 264/265 et la détection des
annulations reste interne à chaque fichier. L'annexe \ref{ann:qualite} documente
huit anomalies ou particularités effectivement présentes avec leurs lignes sources.
